%%%%%%%%%%%%%%%%%%%%%%%%%%%%%%%%%%%%%%%%%%%%%%%%%%%%%%%%%%%%%%%%%%%%%%%%%%%%%%%%
%  01-SAMPUL.TEX — HALAMAN JUDUL (SLIDE 1)
%  -----------------------------------------------------------------------------
%  Slide ini otomatis menampilkan data dari konfigurasi.tex (\title, \author,
%  \institute, \date) melalui perintah \titlepage — tidak perlu mengedit apa
%  pun di sini, kecuali ingin mengubah latar atau susunan logo.
%
%  [plain] = slide tanpa bilah navigasi atas & tanpa footer (bersih untuk
%  halaman judul). Logo pojok otomatis ikut tersembunyi di frame plain.
%%%%%%%%%%%%%%%%%%%%%%%%%%%%%%%%%%%%%%%%%%%%%%%%%%%%%%%%%%%%%%%%%%%%%%%%%%%%%%%%

% \section{} kosong: sampul tidak dihitung sebagai bab di navigasi/daftar isi
\section{}

%--- LATAR SAMPUL --------------------------------------------------------------
% telkom-kampus.jpg dipakai sebagai latar dengan transparansi 20% agar teks
% judul tetap terbaca. Ganti gambarnya di folder assets/ bila perlu.
\usebackgroundtemplate{\transparent{0.2}\includegraphics[width=\paperwidth]{telkom-kampus.jpg}}

\begin{frame}[plain]
  \vspace{-0.5cm}
  \titlepage  % <- otomatis: judul + penulis + institusi + tanggal (dari konfigurasi.tex)

  %--- TIGA LOGO BERJAJAR DI ATAS SLIDE ----------------------------------------
  % Posisi dinyatakan dalam koordinat cm dari tengah slide (0,0):
  %   (3.5, 6.2) = logo kiri | (7.5, 6.2) = tengah | (11.5, 6.2) = kanan
  % Ganti gambar/ukuran sesuai kebutuhan institusi Anda.
  \begin{tikzpicture}[overlay, remember picture]
    \node[inner sep=0pt] at (3.5, 6.2)  {
\includegraphics[width=2.5cm]{logo-telkom-horizontal.png}};
    \node[inner sep=0pt] at (7.5, 6.2)  {
\includegraphics[width=3.5cm]{logo-aicoms.png}};
    \node[inner sep=0pt] at (11.5, 6.2) {
\includegraphics[width=2.5cm]{logo-telkom-horizontal.png}};
  \end{tikzpicture}
\end{frame}

%--- KEMBALIKAN TAMPILAN NORMAL UNTUK SLIDE BERIKUTNYA -------------------------
% Hapus latar gambar; logo pojok kanan atas & footer muncul kembali otomatis.
\usebackgroundtemplate{}
